\documentclass[11pt]{article}
\usepackage{amsmath}
\usepackage{amssymb}
\usepackage{fontspec}
\usepackage{unicode-math}
\setmainfont{DejaVu Serif}
\setsansfont{DejaVu Sans}
\setmonofont{DejaVu Sans Mono}
\setmathfont{DejaVu Math TeX Gyre}
\usepackage[a4paper,margin=2.5cm]{geometry}
\usepackage{booktabs}
\usepackage{longtable}
\usepackage{array}
\usepackage{enumitem}
\usepackage{xcolor}
\usepackage[colorlinks=true,linkcolor=blue!50!black,urlcolor=blue!50!black,citecolor=blue!50!black]{hyperref}
\usepackage{parskip}
\usepackage{fancyhdr}
\pagestyle{fancy}
\fancyhf{}
\renewcommand{\headrulewidth}{0.4pt}
\fancyhead[C]{\footnotesize ARob -- UAVs / Aeronaves Robotizadas}
\fancyfoot[C]{\thepage}
\setlength{\headheight}{26pt}
\sloppy
\emergencystretch=2em
\title{Lab 2 - Estimation of Motion Variables}
\author{Course notes -- 2025/26 labs}
\date{}
\begin{document}
\maketitle
\tableofcontents
\vspace{1em}

Source files:
\begin{itemize}
\item \texttt{25\_\allowbreak{}26\_\allowbreak{}UAVs\_\allowbreak{}Lab2.pdf} (6 pages) - lab handout, "Estimation of Motion Variables of the Parrot AR.Drone," Técnico Lisboa, MSc Aerospace Engineering, 2025/2026 First Semester, October 2025.
\item \texttt{Partial attitude and rate gyro bias estimation observability analysis filter design and performance evaluation.pdf} (9 pages) - P. Batista, C. Silvestre, P. Oliveira, \textit{International Journal of Control}, vol. 84, no. 5, pp. 895-903, May 2011.
\item \texttt{DevKit\_\allowbreak{}Lab2/\allowbreak{}ARDroneSimulinkDevKit\_\allowbreak{}NAV/} - all \texttt{.m} scripts (\texttt{start\_\allowbreak{}here\_\allowbreak{}NAV.m}, \texttt{lib/\allowbreak{}decode\_\allowbreak{}32bit.m}, \texttt{lib/\allowbreak{}getWaypoints.m}, \texttt{lib/\allowbreak{}setupARModel.m}, \texttt{Replay/\allowbreak{}setupReplay.m}, \texttt{simulation/\allowbreak{}setupBaseModel.m}, \texttt{simulation/\allowbreak{}setupHoverSim.m}, \texttt{simulation/\allowbreak{}setupWPTrackingSim.m}, \texttt{wifiControl/\allowbreak{}setupHover.m}, \texttt{wifiControl/\allowbreak{}setupWPTracking.m}).
\end{itemize}

Handout is in English. All pages/files processed in full.

\section{Lab handout - objectives and tasks}

\textbf{Objective}: design, implement, and analyze several estimation solutions for the \textbf{roll and pitch angles} of the Parrot AR.Drone quadrotor, in simulation and experimentally.

\textbf{Deliverable}: single-column PDF report, max 15 pages, plus the \texttt{.m} scripts, submitted via Fenix using the official cover page (\texttt{25\_\allowbreak{}26\_\allowbreak{}UAVs\_\allowbreak{}Lab\_\allowbreak{}report\_\allowbreak{}cover\_\allowbreak{}page.docx}, already in this project). Questions are tagged \textbf{(T)} theoretical (to be solved before the lab session) or \textbf{(L)} laboratory (simulation/experiment, done during the session).

\textbf{Section 2 - Setup and experiments.} New devkit features vs. Lab 1:
\begin{itemize}
\item Entry script is now \texttt{start\_\allowbreak{}here\_\allowbreak{}NAV.m} (adds a 3rd top-level option: replay stored data).
\item Wi-Fi control modes now log raw navigation data into a workspace variable \texttt{navdata} at \textbf{200 Hz}.
\item The navdata-decoding block is an "extended version" exposing \textbf{accelerometer and rate-gyro} samples (Lab 1's devkit apparently only exposed velocity/attitude estimates, not raw IMU data).
\end{itemize}

Five data-collection experiments define the whole lab's empirical basis:
\begin{itemize}
\item \textbf{Experiment A}: sensors at rest, motors off, Hover mode, no take-off, ~30 s - baseline/quiescent noise characterization.
\item \textbf{Experiment B}: on the ground, motors spinning, hull held down to prevent lift-off, take-off command sent, 30 s, then land - isolates motor-vibration-induced noise.
\item \textbf{Experiment C}: real hover flight, take-off → 30 s hold → land - full operational noise (motor vibration + aerodynamic disturbance + control activity).
\item \textbf{Experiment D}: step response in pitch reference, by editing \texttt{ARDroneHover.slx} (same style of step test as in Lab 1) - used to see filter behavior under a real, non-stationary attitude change.
\item \textbf{Experiment E}: short waypoint mission - edit \texttt{ARDroneHover.slx} to give a 2-waypoint reference in $y$: origin → $[0,0,0.75]$ → $[0,1,0.75]$ → back to $[0,0,0.75]$ → land, 5 s dwell between reference changes. Used to see estimator behavior across a realistic short mission.
\end{itemize}

\textbf{Section 3 - Modeling and characterization of sensors} (all "(L)"):
\begin{itemize}
\item 3.1-3.3: compute mean and covariance of accelerometer + rate-gyro readings for Experiments A, B, C respectively, and discuss \textit{why} uncertainty grows from A→B→C (new disturbance sources: motor vibration in B, aerodynamic/control-induced disturbance in C).
\item 3.4: compute "raw" pitch/roll from the accelerometer alone (the inclinometer formulas $\hat\phi_{accel} = \operatorname{atan2}(...)$, $\hat\theta_{accel} = \sin^{-1}(...)$ documented in \href{../../lectures-22-23/notes/03-sensors-estimation.pdf}{\texttt{../../lectures-22-23/notes/03-\allowbreak{}sensors-\allowbreak{}estimation.pdf}}), to serve as the \textit{baseline} against which every filter built later in the lab is compared.
\end{itemize}

\textbf{Section 4 - Kalman filters} (pitch first, then roll by symmetry in 4.7):
\begin{itemize}
\item 4.1: design/implement a \textbf{steady-state Kalman filter} for pitch using the pitch inclinometer measurement $\theta_m$ and rate-gyro measurement $w_{ym}$ (note: MATLAB's \texttt{kalman()} or \texttt{estim()} commands are explicitly suggested - i.e. this is the LTI/steady-state KF machinery from \href{../../lectures-22-23/notes/02-control.pdf}{\texttt{../../lectures-22-23/notes/02-\allowbreak{}control.pdf}} Ch4, not a hand-rolled recursive filter).
\item 4.2: derive, for the constant (steady-state) Kalman gains, the two closed-form transfer functions $\theta_m \to \hat\theta$ and $w_{ym} \to \hat\theta$ - this is exactly the \textbf{complementary filter transfer-function pair} $G_L(s)$/$G_H(s)$ documented in the lecture notes, just derived here from the KF gain rather than assumed a priori.
\item 4.3: discuss effect of tuning $Q$ (process noise weight) vs $R$ (measurement noise weight) - same Bryson's-rule-style tradeoff as Ch4's LQR tuning discussion, applied here to the dual (estimation) problem.
\item 4.4: \textbf{augment the state with a gyro-bias state} $b_y$, model $\dot b_y = 0$ and $w_{ym} = q + b_y$ (bias-augmented rate-gyro model). Suggested validation: inject an artificial bias in simulation and plot true vs. estimated bias to tune the filter. This is a \textbf{2-state, 2-output (or 2-state, single-channel with 2 measurement paths) special case} of exactly the same architecture as the full 6-state paper filter in Section 5, but restricted to a single axis (pitch, then roll) instead of jointly estimating a 3-vector attitude + 3-vector bias.
\item 4.5: check whether the bias-free and bias-augmented filters are \textbf{complementary filters} in the formal sense ($G_L+G_H=1$), and discuss the tradeoffs of that filter class (fast high-frequency response from the gyro, but drift if not bias-compensated; slow but bias-free low-frequency truth from the inclinometer).
\item 4.6: relate results back to Experiments A-E (i.e. show how filter performance/steady-state error changes with the noise regime of each experiment).
\item 4.7: repeat everything for \textbf{roll}, using the roll inclinometer and rate-gyro measurement $w_{xm}$ (by symmetry with pitch/$w_{ym}$).
\end{itemize}

\textbf{Section 5 - Integrated roll and pitch estimation} (the capstone task, "(L)" only, no separate theory sub-parts spelled out): implement the \textbf{Kalman filter for the "nominal system (2)"} of Batista/Silvestre/Oliveira (2011) - i.e. NOT the full nonlinear/GAS-error-dynamics observer from the paper's Theorem 3.4, but the simpler linear time-varying system (their Eq. 2) treated as a standard continuous-time Kalman filter (their Section 4, "Filter design"). The handout explicitly narrows scope: \textbf{only the accelerometer-based partial-attitude vector observation + 3-axis rate-gyro-bias estimation} is required, i.e. the exact 6-state filter of the paper ($x = [y; b_\omega] \in \mathbb{R}^6$), not a re-derivation of the observability proofs.

\section{Reference paper - Partial attitude and rate gyro bias estimation (Batista, Silvestre, Oliveira 2011)}

\textbf{Problem}: with only a \textbf{single} body-fixed vector observation of a known constant inertial vector (e.g. gravity via the accelerometer), full attitude ($R \in SO(3)$) cannot be recovered - one rotational degree of freedom (rotation about the observed vector, i.e. yaw when the vector is gravity) is fundamentally unobservable, since $y(t) = R^T(t)\, {}^Iy = R^T(t)R_0^T\, {}^Iy$ for any rotation $R_0$ about ${}^Iy$. The paper's contribution: exploit \textbf{biased rate-gyro measurements} in addition to the single vector observation to (a) filter the vector observation itself, and (b) fully estimate the 3-axis gyro bias, with a filter that has \textbf{globally asymptotically stable (GAS)} error dynamics - stronger than typical EKF-based attitude filters, which only have local/linearization-based guarantees.

\textbf{System model} (their Eq. 1-2): let $y(t) \in \mathbb{R}^3$ be the vector observation in body coordinates (${}^Iy = R(t)y(t)$, ${}^Iy$ constant, known, in the inertial frame), $\omega_m(t) = \omega(t) + b_\omega(t)$ the biased rate-gyro reading, $b_\omega$ the (assumed constant) bias. Using $\dot y = -S(\omega)y$ (differentiating the observation equation, since $\dot R = RS(\omega)$) and substituting the biased measurement:
\begin{verbatim}
ẋ1(t) = -S(ωm(t))x1(t) + S(x2(t))x1(t)     [x1 = y, using ω = ωm - bω and x×y=-y×x]
ẋ2(t) = 0                                    [x2 = bω, constant]
y(t)  = x1(t)
\end{verbatim}
Compact linear time-varying (LTV) form - \textbf{this is the crux of the whole paper}: although the system is \textit{physically} nonlinear (the state $x_2$ multiplies $x_1$ in the dynamics), it can be rewritten so that the \textit{output} $y(t)$ appears in the system matrix instead of the state, making it \textbf{linear time-varying for observer-design purposes} (the time-variation being driven by the measured/known signals $\omega_m(t)$ and $y(t)$, not by the unknown state):
\[
\dot x(t) = A(t)x(t),\quad y(t) = Cx(t),\quad x = [x_1; x_2] \in \mathbb{R}^6,
\]
\[
A(t) = \begin{bmatrix}-S(\omega_m(t)) & -S(y(t)) \\ 0 & 0\end{bmatrix} \in \mathbb{R}^{6\times 6},\quad C = [I,\ 0] \in \mathbb{R}^{3\times 6}.
\]
This is the \textbf{"nominal system (2)"} the lab handout explicitly asks students to implement a Kalman filter for.

\textbf{Observability analysis} (Section 3, the paper's main theoretical contribution):
\begin{itemize}
\item Lemma 3.1: for a broader class of nonlinear systems expressible this way, observability of the underlying nonlinear system follows from invertibility of the \textbf{observability Gramian} of the associated LTV pair.
\item Theorem 3.3: system (2) is observable on $[t_0,t_f]$ iff for every unit vector $d \in \mathbb{R}^3$, there exists $t^* \in [t_0,t_f]$ with $\|y(t^*)\times d\| > 0$ - i.e., \textbf{the vector observation $y(t)$ must not stay constant/parallel to any fixed direction throughout the interval} (the vehicle must actually rotate/maneuver relative to the observed vector). A vehicle sitting motionless in hover, with $y(t)$ = constant gravity direction in body frame, would make this system \textbf{unobservable} - directly analogous to, and a strict generalization of, the course's own "rate-gyro-only → not observable" worked example in \href{../../lectures-22-23/notes/03-sensors-estimation.pdf}{\texttt{../../lectures-22-23/notes/03-\allowbreak{}sensors-\allowbreak{}estimation.pdf}} (there, a \textit{constant-turn} roll estimator failed observability from yaw-rate alone; here, a \textit{stationary} attitude/bias estimator fails observability from a single constant vector reading alone - both need an added, independent source of information/persistent excitation to become observable).
\item Theorem 3.4 (stronger, needed for GAS filter design): \textbf{uniform complete observability} holds iff there exist $\alpha,\delta>0$ such that for every unit $d$ and every $t$, some $t^* \in [t,t+\delta]$ gives $\|y(t^*)\times d\| \geq \alpha$ - a \textbf{persistency-of-excitation} condition (not just "eventually observable," but observable with a uniform margin on every sliding window of length $\delta$). This is the condition that actually gets invoked to prove the Kalman filter below has GAS error dynamics (via Kalman \& Bucy 1961 / Jazwinski 1970 general stability results for uniformly completely observable LTV Kalman filters).
\end{itemize}

\textbf{Filter design (Section 4)}: because (2) is treated as LTV (not nonlinear) for filtering purposes, the \textbf{standard continuous-time Kalman filter} applies directly, with additive process/measurement noise $w(t)$, $n(t)$ (covariances $\Xi$, $\Theta$):
\[
\dot{\hat x}(t) = A(t)\hat x(t) + K(t)[y(t) - C\hat x(t)],
\]
\[
\dot P(t) = A(t)P(t) + P(t)A(t)^T + \Xi - P(t)C^T\Theta^{-1}CP(t) \quad \text{(Riccati differential equation)},
\]
\[
K(t) = P(t)C^T\Theta^{-1}.
\]
Because $(A(t),C)$ is uniformly completely observable (Theorem 3.4), the filter error dynamics are \textbf{globally asymptotically stable} - this is the paper's headline result, and it is the reason the course picked this filter as the "advanced" capstone estimator for Lab 2 rather than a generic EKF.

\textbf{Simulation results (Section 5)}: accelerometer (gravity vector) + biased rate gyro; injected bias $b_\omega=[2,-3,1]^T\cdot\pi/180$ rad/s; sensor noise std 0.05 m/s² (accel) and 0.05°/s (gyro); $\Xi=\operatorname{diag}(0.05I,0.01I)$, $\Theta=0.05I$. Recovers roll and pitch (the two Euler angles obtainable from a single vector = gravity) with small residual error (~tenths of a degree) and converges the gyro-bias estimate within ~20 s.

\textbf{Experimental results (Section 6)}: validated against a precision motion rate table (Ideal Aerosmith 2103HT, 0.00025° resolution ground truth) with a low-cost MEMSENSE NANO IMU (accel std 0.008 m/s², gyro std 0.95°/s, 150 Hz) - i.e. sensor quality comparable to what the AR.Drone's onboard IMU is expected to have. Filter converges fast ($<1$ s) even from a zero-bias initial guess, and steady-state roll/pitch error stays within about $\pm 0.5$-$1°$ despite the raw accelerometer-only measurement being noisy ($\sim\pm 1$-$2°$) - a concrete, quantified demonstration of the filtering benefit that Lab 2's Sections 3-4 ask students to reproduce qualitatively on the AR.Drone.

\textbf{Important note for the lab report}: the paper's $x_2 = b_\omega$ is the \textit{full 3-axis} gyro bias and $x_1 = y$ is the (3-component) vector observation itself (not roll/pitch angles directly) - recovering roll and pitch from the filtered $\hat x_1$ requires the same $\operatorname{atan2}$/$\sin^{-1}$ inversion documented in \href{../../lectures-22-23/notes/03-sensors-estimation.pdf}{\texttt{../../lectures-22-23/notes/03-\allowbreak{}sensors-\allowbreak{}estimation.pdf}}'s "sensor-model inversion" section, applied post-filter instead of pre-filter (i.e. filter the raw vector first, then convert to angles - this generally gives better results than converting to angles first and filtering the angles, since the vector-space dynamics are the ones that are exactly LTV/exploitable, while angle-space dynamics involve trig nonlinearities).

\section{Devkit guide/scripts - NAV-specific additions vs Lab 1's devkit}

\begin{center}
\small
\setlength{\tabcolsep}{3pt}
\begin{longtable}{|p{3.0cm}|p{4.2cm}|p{7.0cm}|}
\hline
Script & Purpose & Notable content \\
\hline
\endhead
\texttt{start\_\allowbreak{}here\_\allowbreak{}NAV.m} & Entry point, replaces Lab 1's \texttt{start\_\allowbreak{}here.m} & Adds top-level \textbf{option (3) Replay from stored data}: prompts for a \texttt{.mat} filename, loads it, \texttt{cd Replay}, runs \texttt{setupReplay}. Options (1) Simulation and (2) Wi-Fi control are structurally identical to Lab 1's devkit (same sub-menus: baseline/hover/waypoint-tracking for simulation; hover/waypoint-tracking for Wi-Fi). \\
\hline
\texttt{lib/\allowbreak{}decode\_\allowbreak{}32bit.m} & Decodes raw 32-bit navdata fields from the AR.Drone's UDP packet stream & Takes a 4-byte input array, reassembles it (with byte-order swap: \texttt{hex\_\allowbreak{}value(4,3,2,1)}), and returns either a \texttt{uint32} (\texttt{mark=1}), a signed value via \texttt{hex2dec(...)-2\textasciicircum31} (\texttt{mark=2}), or a \texttt{single}-precision float reinterpretation of the same 4 bytes (\texttt{mark} otherwise) via \texttt{typecast}. This is the low-level function that turns the drone's raw binary telemetry (battery status, accelerometer/gyro raw counts, etc.) into usable MATLAB numeric values - genuinely new vs. Lab 1, which per the handout only exposed higher-level velocity/attitude estimates, not raw IMU counts. \\
\hline
\texttt{lib/\allowbreak{}getWaypoints.m} & Same as Lab 1 - hardcoded 5-column waypoint list $[X_e, Y_e, h, \psi, \text{wait\_\allowbreak{}time}]$, edited per-mission (currently a small square path, take off → 2 corners → back → land, 5 s dwell each) & Unchanged in structure from Lab1; a large commented-out block shows an "ORIGINAL" 10-waypoint list, suggesting the file is meant to be edited per-experiment (e.g. for Experiment E's waypoint mission). \\
\hline
\texttt{lib/\allowbreak{}setupARModel.m} & Loads linearized transfer-function / state-space models (\texttt{ssRoll}, \texttt{ssRoll2V}, \texttt{ssPitch}, \texttt{ssPitch2U}, \texttt{ssYaw}, \texttt{ssH}) from \texttt{.mat} files & Identical to Lab 1's devkit; these are the small-signal linear models used by the \textit{simulation} path only (not by the raw-navdata-based estimators built in this lab, which work directly on \texttt{navdata}). \\
\hline
\texttt{Replay/\allowbreak{}setupReplay.m} & Replay entry point & Just adds \texttt{../lib} to path and opens \texttt{ARDroneReplay\_\allowbreak{}V2.slx} - the loaded \texttt{.mat} (containing a previously-recorded \texttt{navdata} variable, per \texttt{start\_\allowbreak{}here\_\allowbreak{}NAV.m}'s \texttt{load(filename)} call) is fed into this model instead of live Wi-Fi data, letting students re-run/develop estimators offline against Experiment A-E recordings without needing the physical drone each time. \\
\hline
\texttt{simulation/\allowbreak{}setupBaseModel.m}, \texttt{setupHoverSim.m}, \texttt{setupWPTrackingSim.m} & Simulation-mode setup scripts & Unchanged from Lab 1's devkit: set \texttt{FMS.Ts=0.065s}, \texttt{timeDelay=4·FMS.Ts}, load the linear model, waypoints, \texttt{simDT=0.005s}, then open the corresponding \texttt{.slx} (\texttt{ARDroneBaseSim}, \texttt{ARDroneHoverSim}, \texttt{ARDroneWPTrackingSim}). \\
\hline
\texttt{wifiControl/\allowbreak{}setupHover.m}, \texttt{setupWPTracking.m} & Wi-Fi control setup scripts & Sample time now explicitly \texttt{1/200} (200 Hz, matching the new navdata logging rate quoted in the handout) rather than the simulation's \texttt{0.065s}; both carry an explicit warning comment that the onboard \textbf{position estimate is inaccurate} because it integrates an already-inaccurate optical-flow-based velocity estimate - a direct, practical illustration of why Lab 2's whole point (building better estimators from raw IMU data) matters. Opens \texttt{ARDroneHover\_\allowbreak{}V2.slx} / \texttt{ARDroneWPTracking\_\allowbreak{}V2.slx} (the "\texttt{\_\allowbreak{}V2}" suffix distinguishing these from Lab 1's plain \texttt{ARDroneHover.slx}/\texttt{ARDroneWPTracking.slx}, which are also still present in this devkit and are what Experiment D's instructions say to edit for the pitch-step test). \\
\hline
\end{longtable}
\end{center}

\section{Simulink models present}

All in \texttt{DevKit\_\allowbreak{}Lab2/\allowbreak{}ARDroneSimulinkDevKit\_\allowbreak{}NAV/}:
\begin{itemize}
\item \texttt{lib/ARBlocks.slx} - core AR.Drone block library (dynamics/interface blocks reused by every other model).
\item \texttt{simulation/\allowbreak{}ARDroneBaseSim.slx}, \texttt{ARDroneHoverSim.slx}, \texttt{ARDroneWPTrackingSim.slx} - pure-simulation models (linear dynamics, no real drone).
\item \texttt{wifiControl/\allowbreak{}ARDroneHover.slx}, \texttt{ARDroneWPTracking.slx} - original (Lab 1-era) Wi-Fi control models; per the handout, these are the ones Experiments D and E instruct students to modify (add a pitch step / a 2-waypoint $y$ reference).
\item \texttt{wifiControl/\allowbreak{}ARDroneHover\_\allowbreak{}V2.slx}, \texttt{ARDroneWPTracking\_\allowbreak{}V2.slx} - updated Wi-Fi control models with the extended navdata decoding (200 Hz raw accel/gyro logging) - these are the ones actually launched by \texttt{setupHover.m}/\texttt{setupWPTracking.m} and used to run Experiments A-C.
\item \texttt{Replay/\allowbreak{}ARDroneReplay\_\allowbreak{}V2.slx} - offline replay model consuming a previously-logged \texttt{navdata} variable instead of a live Wi-Fi feed.
\end{itemize}

\section{Connections to lecture material}

\begin{itemize}
\item \textbf{Raw pitch/roll from accelerometer (handout 3.4)} $\leftrightarrow$ the exact inclinometer formulas $\hat\phi_{accel} = \operatorname{atan2}(y_{accel,y}, y_{accel,z})$-style, $\hat\theta_{accel} = \sin^{-1}(y_{accel,x}/g)$ derived in \href{../../lectures-22-23/notes/03-sensors-estimation.pdf}{\texttt{../../lectures-22-23/notes/03-\allowbreak{}sensors-\allowbreak{}estimation.pdf}} under "Sensor-model inversion," including the explicit caveat there that this approach fails during maneuvers (motivating everything that follows in both the lecture and the lab).
\item \textbf{Steady-state KF for pitch (4.1-4.3)} $\leftrightarrow$ directly the LTI Kalman filter machinery from \href{../../lectures-22-23/notes/02-control.pdf}{\texttt{../../lectures-22-23/notes/02-\allowbreak{}control.pdf}} (Ch4: \texttt{K=lqr(...)}-style tooling, duality with \texttt{L=place(...)}) and the continuous-discrete KF equations in \texttt{03-\allowbreak{}sensors-\allowbreak{}estimation.pdf} ($A_d=I+ATs+A^2Ts^2/2$, $P_{k+1}=A_dP_kA_d^T+Ts^2Q$), specialized here to MATLAB's built-in \texttt{kalman()}/\texttt{estim()} for the steady-state (algebraic Riccati) solution rather than a manually-coded recursive filter.
\item \textbf{Bias-augmented KF (4.4)} $\leftrightarrow$ the exact "second-order complementary filter" structure documented in \texttt{03-\allowbreak{}sensors-\allowbreak{}estimation.pdf} ($x=[\phi,b_p]$, $u=p_m$, $y=\phi_m$, giving $\hat\Phi(s)=G_L(s)\Phi_m(s)+G_H(s)P_m(s)/s$) - the lab has students re-derive/re-tune this same structure for pitch and roll on real AR.Drone data.
\item \textbf{Complementarity check (4.5)} $\leftrightarrow$ \texttt{03-\allowbreak{}sensors-\allowbreak{}estimation.pdf}'s explicit statement that "the complementary filter is a Kalman filter with fixed (non-adaptive) gain" - Section 4.5 of the lab is literally asking students to verify this equivalence ($G_L+G_H=1$) for their own designed filters.
\item \textbf{Section 5's paper-based filter} $\leftrightarrow$ generalizes the course's own three-stage worked example in \texttt{03-\allowbreak{}sensors-\allowbreak{}estimation.pdf} ("rate-gyro only → not observable," "+ accelerometer → observable"): the Batista/Silvestre/Oliveira filter is the natural full-3D, full-rigor version of that same idea - single vector observation (accelerometer/gravity) + biased rate gyro, formalized with a proper observability proof and a GAS convergence guarantee, instead of the course's simplified scalar/coordinated-turn special case.
\item The \textbf{LQR $\leftrightarrow$ Kalman filter duality} explicitly taught in Ch4/Ch6 ($K=R^{-1}B^TP$ vs $L=\Sigma C^TN^{-1}$, dual Riccati equations) is exactly the mathematical machinery underlying the paper's Riccati equation $\dot P=AP+PA^T+\Xi-PC^T\Theta^{-1}CP$ and gain $K=PC^T\Theta^{-1}$ - same structure, just continuous-time LTV instead of LTI.
\item The \textbf{AR.Drone-specific note in \texttt{03-\allowbreak{}sensors-\allowbreak{}estimation.pdf}} that the AR.Drone has no GPS and instead uses sonar + downward-camera vision is corroborated here: \texttt{setupHover.m}/\texttt{setupWPTracking.m}'s explicit in-code warning that the onboard position estimate (integrated optical-flow velocity) is inaccurate is precisely why Lab 2 focuses only on attitude (roll/pitch) estimation from IMU data, not full position/velocity estimation.
\end{itemize}

\section{Open questions / unclear points}

\begin{itemize}
\item The handout's Eq. numbering ("system (2)") refers to the \textit{paper's} equation 2, not a lab-internal numbering - confirmed by cross-checking against the paper itself (its Eq. (2) is exactly the compact LTV form $\dot x=A(t)x$, $y=Cx$ derived above).
\item The handout does not specify numeric noise covariances ($Q$, $R$ / $\Xi$, $\Theta$) to use for the AR.Drone - students are expected to derive these from the Experiment A-C mean/covariance characterization in Section 3, unlike the reference paper which states its simulation/experimental covariances explicitly. Matching the AR.Drone's real IMU noise characteristics to the paper's simulated/experimental values (from a much higher-grade MEMSENSE IMU) is left entirely to the student.
\item \texttt{decode\_\allowbreak{}32bit.m}'s \texttt{mark=2} branch (\texttt{hex2dec(hex\_\allowbreak{}value)-2\textasciicircum31}) is used for signed integers but the exact set of navdata fields that require each of the three \texttt{mark} modes (unsigned int / signed int / float) is not documented in the script itself - would need the Parrot AR.Drone SDK's navdata structure reference to map specific byte offsets to sensor channels.
\item Which specific navdata fields (byte offsets into the \texttt{NDS} structure referenced in \texttt{decode\_\allowbreak{}32bit.m}'s comment, e.g. accelerometer/gyro sample locations) are wired into the "extended version" navdata-decoding Simulink block is not visible from the \texttt{.m} files alone (the block internals are in the binary \texttt{ARBlocks.slx} / \texttt{ARDroneHover\_\allowbreak{}V2.slx}, not inspected here) - would require opening these \texttt{.slx} files in Simulink to confirm exact scaling/units of the raw accelerometer and rate-gyro channels before Section 3's mean/covariance characterization can be done correctly.
\item The paper's filter is explicitly noted (Remark 4) to not be truly "optimal" in a strict sense because of multiplicative noise (the $S(y(t))$ and $S(\omega_m(t))$ terms depend on noisy measured/estimated signals, not clean known signals) - the lab handout does not ask students to address this caveat, but it's worth flagging if a report reviewer probes filter optimality claims.
\end{itemize}

\end{document}
